%! Roots and Polynomials — Unit 3, Lesson 3.6 — AP Practice (ANSWER KEY)
% EXTRA CREDIT — optional, scored 1 to 5, exactly TWO pages, last in the packet (graded homework
% is DeltaMath). Page 1 is Section I, five multiple-choice items with four options (A)–(D); page 2
% is Section II, one multi-part free-response question on a table. Its contexts (a sound
% engineer's polynomial, an unlabeled cubic graph h, a roller-coaster track H) are used nowhere
% else in the lesson.
% Section II has three parts, each with a 2in rubric box (math-shared's minimum), so page 2 holds
% exactly three. Distractors are the lesson's errors on purpose: MC 1 flips the factor signs (A, D) or drops the
% factor x for the zero 0 (C); MC 2 negates the wrong part of 4 - i (A, C) or drops the imaginary
% part (B); MC 3 forgets the conjugate (B), counts only real zeros (A), or conjugates the real
% zeros too (D); MC 4 flips the factor signs (B), forgets a (C), or swaps the multiplicities (D);
% MC 5 inverts a (A), takes the y-intercept as a (D), or flips a factor sign (C).
%
% PAGE LOCKSTEP with the key: the key marks the correct option in its LABEL only, and every
% \gradedblank line is identical in both files. The explicit \newpage fixes the page break.
\documentclass[10pt]{article}
\usepackage{precalculus-article}
\usepackage{precalculus-key}
\usepackage{grading}
\gradingassignment{precalc-u03-l06-ap}
\gradingcourse{precalculus}
\newcommand{\TallMath}[1]{$\displaystyle #1\rule[-1.4em]{0pt}{3.2em}$}

\begin{document}

\pageheader{Unit 3, Lesson 3.6}{AP Practice --- Answer Key}
{\small\itshape Write each final answer in its box. Only what is inside a box is graded.}\par

\vspace{0.08in}

\boxguard[8]
\begin{remindbox}
\small
\textbf{Extra credit, optional, scored 1 to 5.} These two pages are written in the style of the AP
Precalculus exam. Your graded homework is on DeltaMath; do it first. Every
answer choice below is a mistake someone really makes, so read all four before you choose.
\end{remindbox}

\vspace{0.1in}

\begin{headlinebox}{goldbg}
\small
\textbf{Section I --- Multiple Choice.} Write the letter of the single best answer in its Letter box. No calculator.
\end{headlinebox}

\vspace{0.06in}

\begin{enumerate}[leftmargin=1.9em, itemsep=12pt]

\item A sound engineer needs a polynomial with leading coefficient 1 whose only zeros are
      $-3$, $0$, and $5$, each with multiplicity 1. Which of the following could it be?

      \smallskip
      \begin{tabular}{@{}p{4.2cm}p{4.2cm}@{}}
      \textbf{(A)} $x(x - 3)(x + 5)$ & {\color{keyred}\textbf{$\checkmark$\,(B)}} $x(x + 3)(x - 5)$ \\[2pt]
      \textbf{(C)} $(x + 3)(x - 5)$ & \textbf{(D)} $(x - 3)(x + 5)$
      \end{tabular}
      \hfill\textit{Letter:}~\gradedblank[id=ap1, type=choice, points=1, answer=B]{0.42in}{13pt}

\item The polynomial $p$ has real coefficients, and $4 - i$ is a zero of $p$. Which of the
      following must also be a zero of $p$?

      \smallskip
      \textbf{(A)} $-4 + i$ \hspace{2.2em} \textbf{(B)} $4$ \hspace{2.2em} \textbf{(C)} $-4 - i$ \hspace{2.2em} {\color{keyred}\textbf{$\checkmark$\,(D)}} $4 + i$ \hfill\textit{Letter:}~\gradedblank[id=ap2, type=choice, points=1, answer=D]{0.42in}{13pt}

\item A polynomial has real coefficients and zeros $2$, $-2$, and $3 + i$. What is the least
      possible degree of the polynomial?

      \smallskip
      \textbf{(A)} $2$ \hspace{2.2em} \textbf{(B)} $3$ \hspace{2.2em} {\color{keyred}\textbf{$\checkmark$\,(C)}} $4$ \hspace{2.2em} \textbf{(D)} $6$ \hfill\textit{Letter:}~\gradedblank[id=ap3, type=choice, points=1, answer=C]{0.42in}{13pt}

\item \begin{minipage}[t]{0.58\linewidth}
      The graph of the cubic polynomial $h$ is shown. It crosses the $x$-axis at $x = -1$,
      touches it at $x = 2$, and passes through $(0, -8)$. Which of the following could
      define $h$?
      \begin{enumerate}[leftmargin=2.6em, itemsep=3pt]
        \item[{\color{keyred}$\checkmark$\,(A)}] $h(x) = -2(x + 1)(x - 2)^2$
        \item[(B)] $h(x) = -2(x - 1)(x + 2)^2$
        \item[(C)] $h(x) = (x + 1)(x - 2)^2$
        \item[(D)] $h(x) = -2(x + 1)^2(x - 2)$
      \end{enumerate}
      \vspace{6pt}
      \hfill\textit{Letter:}~\gradedblank[id=ap4, type=choice, points=1, answer=A]{0.42in}{13pt}
      \end{minipage}\hfill
      \begin{minipage}[t]{0.36\linewidth}
      \vspace{-2pt}
      \centering
      \begin{tikzpicture}[x=0.6cm, y=0.15cm, baseline=(current bounding box.north)]
        \draw[linegray, very thin, xstep=1, ystep=2] (-2,-10) grid (4,8);
        \draw[->, thick] (-2,0) -- (4.4,0) node[right] {\scriptsize $x$};
        \draw[->, thick] (0,-10) -- (0,8.8) node[above] {\scriptsize $h(x)$};
        \foreach \x in {-1,1,2,3} \node[below, font=\tiny] at (\x,0) {\x};
        \foreach \y in {-8,-4,4} \node[left, font=\tiny] at (0,\y) {\y};
        \draw[very thick, plum, domain=-1.3:2.95, samples=80, smooth]
             plot (\x, {-2*(\x+1)*(\x-2)*(\x-2)});
        \fill[goldacc] (0,-8) circle (2.2pt) node[right, font=\tiny, fill=white, inner sep=1pt] {$(0, -8)$};
      \end{tikzpicture}
      \end{minipage}

\item The polynomial $q$ has degree 3, zeros $-1$, $2$, and $5$, and $q(0) = 20$. What is the
      leading coefficient of $q$?

      \smallskip
      \textbf{(A)} $\tfrac{1}{2}$ \hspace{2.2em} {\color{keyred}\textbf{$\checkmark$\,(B)}} $2$ \hspace{2.2em} \textbf{(C)} $-2$ \hspace{2.2em} \textbf{(D)} $20$ \hfill\textit{Letter:}~\gradedblank[id=ap5, type=choice, points=1, answer=B]{0.42in}{13pt}

\end{enumerate}

\newpage

\begin{headlinebox}{goldbg}
\small
\textbf{Section II --- Free Response.} Show your reasoning outside the boxes. Answers about the
track must be written \emph{in context}, with units --- that is what earns credit on the AP exam.
Write each final answer in its box. Only what is inside a box is graded.
\end{headlinebox}

\vspace{0.06in}

\begin{enumerate}[leftmargin=1.9em, itemsep=8pt, start=6]

\item A roller-coaster designer models one section of track with a polynomial $H$ of degree 3
      with real coefficients. For $0 \le x \le 6$, $H(x)$ is the height of the track, in meters,
      relative to ground level, where $x$ is the horizontal distance from the station, in tens
      of meters. A negative height means the track is below ground, in a tunnel.

      \begin{center}
      \renewcommand{\arraystretch}{1.35}
      \begin{tabular}{|c|c|c|c|c|c|c|c|}
      \hline
      $x$ (tens of meters) & 0 & 1 & 2 & 3 & 4 & 5 & 6 \\ \hline
      $H(x)$ (meters)      & $-8$ & 0 & 2 & 1 & 0 & 2 & 10 \\ \hline
      \end{tabular}
      \end{center}

      \textbf{(a)} The zeros of $H$ shown in the table are at $x = 1$ and $x = 4$. Use the table
      to explain why the zero at $x = 4$ has even multiplicity. Then write $H(x)$ in factored
      form, using a value from the table to find the leading coefficient.
      \rubricbox{\gradedblank[id=ap6a, type=rubric, points=2,
        display={$H(3) = 1$ and $H(5) = 2$ are both positive, so $H$ does not change sign at $x = 4$: the track touches ground level and rises again, so the multiplicity is even (2). $H(x) = a(x - 1)(x - 4)^2$, and $H(0) = a(-1)(16) = -16a = -8$, so $a = \tfrac{1}{2}$: $H(x) = \tfrac{1}{2}(x - 1)(x - 4)^2$.}]{\linewidth}{\dimexpr 2in-4pt\relax}}

      \textbf{(b)} Write $H(x)$ in standard form. Which form, factored or standard, shows the value
      $H(0)$ without any computing? Interpret $H(0) = -8$ in context, with units.
      \rubricbox{\gradedblank[id=ap6b, type=rubric, points=2,
        display={$\tfrac{1}{2}(x - 1)(x^2 - 8x + 16) = \tfrac{1}{2}x^3 - \tfrac{9}{2}x^2 + 12x - 8$. Standard form shows it: the constant term $-8$ is $H(0)$. At the station ($x = 0$) the track is 8 meters below ground level, in the tunnel.}]{\linewidth}{\dimexpr 2in-4pt\relax}}

      \textbf{(c)} A redesign $G$ must be a cubic with real coefficients and leading coefficient
      $\tfrac{1}{2}$ that keeps the zero at $x = 1$ and has $4 + i$ as a zero. Write $G(x)$, and
      explain why the redesigned track meets ground level only once.
      \rubricbox{\gradedblank[id=ap6c, type=rubric, points=2,
        display={Real coefficients force $4 - i$ too: $G(x) = \tfrac{1}{2}(x - 1)\bigl(x - (4 + i)\bigr)\bigl(x - (4 - i)\bigr) = \tfrac{1}{2}(x - 1)(x^2 - 8x + 17)$. A cubic has exactly 3 zeros; two are the non-real pair $4 \pm i$, so $x = 1$ is the only real zero.}]{\linewidth}{\dimexpr 2in-4pt\relax}}

\end{enumerate}

\end{document}
